\documentclass[10pt,a4paper]{amsart}
\usepackage[margin=19mm]{geometry}
\usepackage{amsmath,amssymb,mathtools}
\usepackage[scr=rsfs,cal=euler]{mathalfa}
\usepackage{xfrac}
\usepackage[unicode,hidelinks,bookmarks=false]{hyperref}
\newcommand{\ie}{i.e.\ }
\newcommand{\cf}{cf.\ }
\newcommand{\resp}{respectively\ }
\newcommand{\Subsection}{\subsection}
\providecommand{\nobreakdash}{\nobreak\hbox{-}}
\setlength{\emergencystretch}{2em}
\multlinegap0pt
\usepackage{mdframed}
\usepackage{mdframed}
\usepackage{mdframed}
\usepackage{mdframed}
\usepackage{braket}
\usepackage{amssymb}
\usepackage[shortlabels]{enumitem}
\usepackage{mathtools}
\usepackage{mdframed}
\newtheorem{proposition}{Proposition}
\title{Cayley's First Hyperdeterminant is an Entanglement Measure}
\author{Isaac Dobes, Naihuan Jing}
\date{Source: arXiv:2504.15511v3 (CC BY 4.0)}
\begin{document}
\maketitle

\noindent\emph{Source and licence.} Cayley's First Hyperdeterminant is an Entanglement Measure. Version arXiv:2504.15511v3 (CC BY 4.0), distributed by arXiv under the Creative Commons Attribution 4.0 International licence, http://creativecommons.org/licenses/by/4.0/, as recorded in the arXiv licence metadata for this exact version and shown on the abstract page https://arxiv.org/abs/2504.15511v3. Full licence notice: \texttt{licenses/arxiv-2504.15511-CC-BY-4.0.txt}.\par\smallskip

\noindent\textit{Editorial scope.} Counted unit: the article's proposition product (source0.tex lines 143--146). The article's complete original proof of it is delivered with it (source0.tex lines 147--160, one \texttt{proof} environment of 14 lines). No mathematics is written, repaired or re-derived: the statement and the proof are the article's own lines, in the article's own notation, and no retained line is substituted or re-flowed.  No further source-local context is needed: every cross-reference of every retained line resolves inside the two delivered spans.  Nothing was omitted from any retained span, so the package carries no omission record.  No retained line contains a citation, so the wrapper carries no bibliography and no bibliography entry.  No retained line contains a figure environment, an image-inclusion command or drawing code, so no image is delivered and the package declares no asset dependency.  Editorial formatting only, in addition: an AMS \texttt{amsart} preamble that restates the article's own declarations and macros used by the retained lines, namely the environments \texttt{proposition} on separate counters, together with the standard packages those lines need.  The retained environments print in this document as Proposition 1 in order of appearance, whereas the article prints the same items with the numbering of its own full document.  The complete native source of the article as distributed in the arXiv e-print for this exact version is bundled unchanged as \texttt{sources/176873/source0.tex}.
\begin{proposition}\label{Kronecker and multi lin mat product}
    Let $\ket{\psi}$ and $\ket{\varphi}$ be pure states of $n$-qudits with corresponding hypermatrix representation $\widehat{\psi}$ and $\widehat{\varphi}$ (respectively). Then for any $d\times d$ matrices $M_1,...,M_n$, we have that  
    \[\ket{\varphi} = (M_1\otimes...\otimes M_n)\ket{\psi} \iff \widehat{\varphi} = (M_1,...,M_n)*\widehat{\psi}.\]
\end{proposition}

\begin{proof} 
    Observe 
    \begin{align*}
        \ket{\varphi} &= (M_1\otimes...\otimes M_n) \ket{\psi} \\
        &= (M_1\otimes...\otimes M_n)\sum\limits_{i_1,...,i_n=0}^{d-1}\psi_{i_1...i_n}\ket{i_1}\otimes...\otimes \ket{i_n} \\
        &= \sum\limits_{i_1,...,i_n=0}^{d-1}\psi_{i_1...i_n}(M_1\ket{i_1})\otimes...\otimes (M_n\ket{i_n}),\quad \text{by linearity of Kronecker products.}
    \end{align*}
    The bijection induced by $H$ maps this to
    \[\sum\limits_{i_1,...,i_n=0}^{d-1}\psi_{i_1...i_n}(M_1\ket{i_1})\circ...\circ (M_n\ket{i_n})\]
    and by the linearity of multilinear matrix multiplication, this is equal to
    \[(M_1,...,M_n)*\left(\sum\limits_{i_1,...,i_n=0}^{d-1}\psi_{i_1...i_n}\ket{i_1}\circ...\circ \ket{i_n}\right)\]
    which is precisely
    \[\widehat{\varphi}=(M_1,...,M_n)*\widehat{\psi}.\]
\end{proof}

\end{document}
